\documentclass[12pt]{article}



\input{macros-sem}

\parindent0pt

%\usepackage{pgf}
%\usepackage{tikz}
%\usepackage[utf8]{inputenc}
%\usetikzlibrary{arrows,automata}
%\usetikzlibrary{positioning}


\newcommand{\probex}{\stepcounter{uebno}{\vspace{0.3cm}\noindent \bf (P\arabic{uebno}) }}
\newcommand{\solution}[1]{$\,$ \\ \begin{proof}#1 \qed \end{proof}}
\newcommand{\solutiong}[1]{$\,$ \\ \begin{proof}\mbox{}\\#1 \qed \end{proof}}

%%%%%%%decomentat in normal, comentat in sol si in sol-grila:
% \renewcommand{\solution}[1]{}
% \renewcommand{\boxtimes}{\Box}

%%%%%%%decomentat in normal si sol, comentat in sol-grila:
% \renewcommand{\solutiong}[1]{}

 
\begin{document}

 
\setcounter{uebno}{0}
\renewcommand{\beschreibung}{Model de examen} 

\renewcommand{\veranstaltung}{{\footnotesize Logic\u a matematic\u a \c si computa\c tional\u a}}

\newsavebox{\tema}
\savebox{\tema}{\raisebox{-11pt}{\parbox[t]{8cm}{\bf {\footnotesize FMI, Info, Anul I}\\ \veranstaltung}}}
\renewcommand{\titel}{ 
\usebox{\tema}  
\bigskip
\vspace{2mm}
\begin{center}
{\Large \bf Model de examen} 


 \end{center} \vspace{5mm}
}

%\pagestyle{empty}

\titel

\thispagestyle{empty}

%\vspace{-1cm}
%\begin{flushleft}
%{
%\textbf{Nume: \underline{\hspace{5cm}}}

%\vspace{.4cm}
%\textbf{Prenume: \underline{\hspace{4cm}}}

%\vspace{.4cm}
%\textbf{Grupa: \underline{\hspace{4.7cm}}}

%}
%\end{flushleft}


%---------------------------------------------
\pagenumbering{arabic}
\vspace{-1cm}
% \begin{flushleft}
% {
% \textbf{Nume: \underline{\hspace{5cm}}}

% \vspace{.4cm}
% \textbf{Prenume: \underline{\hspace{4.4cm}}}

% \vspace{.4cm}
% \textbf{Grupa: \underline{\hspace{4.9cm}}}

% }
% \end{flushleft}

\vspace{.2cm}

\newcommand{\ntaut}{\not\models}
\section*{Partea I. Probleme cu rezolvare clasic\u a}

\probex[1 punct]
Fie $LP$ logica propozi\c tional\u a, $\cE$ o mul\c time de evalu\u ari \c si
$$\Gamma=\{\psi \in Form \mid e\taut \psi \text{ pentru orice } e\in \cE\}.$$
 Presupunem c\u a $\cE$ are cel pu\c tin dou\u a elemente.  Demonstra\c ti c\u a exist\u a o formul\u a $\vp$  cu proprietatea c\u a $\vp\notin \Gamma$ \c si $\neg \vp\notin \Gamma$. 
 
\solution{Fie $e,f$ dou\u a evalu\u ari distincte din $\cE$. Atunci exist\u a $v\in V$ \ai $e(v)\ne f(v)$. Avem dou\u a cazuri:
\be
\item $e(v)=0$ \c si $f(v)=1$, deci $f(\neg v)=0$. Deoarece $e\ntaut v$,  avem c\u a $v\notin \Gamma$. 
Deoarece $f\ntaut \neg  v$, avem c\u a $\neg v\notin \Gamma$.
\item $e(v)=1$ \c si $f(v)=0$. Similar.
\ee
}
\probex[1 punct]
Fie $LP$ logica propozi\c tional\u a \c si $\vp$, $\psi \in Form$. Demonstra\c ti sintactic c\u a:
$$\vdash \vp \si \psi \dnd \psi \si \vp.$$
\solution{ 
Avem:\\[1mm]
\bt{lll}
(1) & $\{\psi,\vp\}\vdash \psi\si \vp$ &  Propozi\c tia 1.68, (48)\\
(2) & $\{\vp,\psi\} \vdash \psi\si \vp$ & deoarece $\{\psi,\vp\}=\{\vp,\psi\}$ \\
(3) & $\{\vp\si \psi\} \vdash \psi\si \vp$ & Propozi\c tia 1.68, (49) \c si (2)\\
(4)  & $\vdash \vp \si \psi \ra \psi \si \vp$ & Teorema deduc\c tiei: (3)\\
(5) & $\vdash  \psi \si \vp \ra \vp \si \psi $ & se demonstreaz\u a similar\\
(6) & $\{ \vp \si \psi \ra \psi \si \vp, \psi \si \vp \ra \vp \si \psi\}\vdash \vp \si \psi \dnd \psi \si \vp$ & Propozi\c tia 1.68, (48) \\
(7) & $\vdash  (\vp \si \psi \ra \psi \si \vp)\ra \big((\psi \si \vp \ra \vp \si \psi)\ra (\vp \si \psi \dnd \psi \si \vp)\big)$ & Teorema deduc\c tiei: (6)\\
(8) & $\vdash (\psi \si \vp \ra \vp \si \psi)\ra (\vp \si \psi \dnd \psi \si \vp)$ & (MP): (4), (7)\\
(9) & $\vdash \vp \si \psi \dnd \psi \si \vp$ & (MP): (5), (8).
\et
}

\probex[1 punct]
Fie $\lc$ un limbaj de ordinul \^int\^ai \c si $\Gamma$ o mul\c time de enun\c turi ale lui $\lc$. 
Demonstra\c ti urm\u atoarele:
\be
\item $Mod(\Gamma)=Mod(Th(\Gamma))$.
\item $Th(\Gamma)$ este o $\lc$-teorie.
\item Fie $T$ o $\lc$-teorie  \ai $\Gamma \se T$. Atunci $Th(\Gamma) \se T$.
\ee
\solution{
\be
\item "$\supseteq$"  Pentru orice $\vp\in \Gamma$, avem c\u a $\Gamma \models \vp$, deci  $\vp\in Th(\Gamma)$.  A\c sadar, $\Gamma\se Th(\Gamma)$.  Aplic\u am  Lema~2.34.(ii). \\
"$\subseteq$" Fie $\cA \in Mod(\Gamma)$ \c si $\vp\in Th(\Gamma)$ arbitrar. Avem c\u a 
$\cA \models \Gamma$ \c si $\Gamma\models \vp$, deci $\cA\models \vp$.  A\c sadar, $\cA\in  
Mod(Th(\Gamma))$.
\item   Demonstr\u am c\u a $Th(\Gamma)$ este o $\lc$-teorie. Pentru orice enun\c t $\vp$,  avem c\u a
\bua 
Th(\Gamma)\models \vp & \Ba & Mod(Th(\Gamma))\se Mod(\vp)
\Ba Mod(\Gamma)\se Mod(\vp) \text{~(conform (i))} \\
& \Ba & \Gamma\models\vp \Ba \vp \in Th(\Gamma).
\eua

\item Fie $T$ o $\lc$-teorie care con\c tine $\Gamma$ \c si $\vp\in Th(\Gamma)$. Din 
$Mod(\Gamma)\se Mod(\vp)$ \c si $Mod(T)\se Mod(\Gamma)$ rezult\u a c\u a $Mod(T)\se Mod(\vp)$, deci
$T\models \vp$. Deoarece $T$ este$\lc$-teorie, ob\c tinem c\u a $\vp\in T$. A\c sadar, $Th(\Gamma)\se T$.
\ee
}

\probex[1 punct]
Fie $\lc$ un limbaj de ordinul \^int\^ai \c si  \c si  $\Gamma$ o mul\c{t}ime de enun\c{t}uri ale lui $\lc$ cu proprietatea c\u{a}
\bce
(*) \quad pentru orice $m\in\N$, $\Gamma$ are un model finit de cardinal  $\geq m$.
\ece
Not\u am cu $\cM$ clasa modelelor finite ale lui $\Gamma$. Demonstra\c ti c\u a $\cM$ nu este axiomatizabil\u{a}.

\solution{Presupunem prin reducere la absurd c\u{a} $\cM$ este axiomatizabil\u{a}. A\c sadar, exist\u{a} $\Delta\se Sen_\lc$  \ai 
\bce
(**) \quad pentru orice $\lc$-structur\u{a} $\cA$, \, $\cA\models \Delta$ \ddacam $\cA\in \cM$.
\ece
Fie 
$$\Sigma:=\Gamma\cup \Delta \cup \{\exists^{\geq n}\mid n\geq 1\}.$$

Demonstr\u{a}m c\u{a} $\Sigma$ este satisfiabil\u{a} folosind Teorema de compacitate. 

Fie $\Sigma_0$ o submul\c{t}ime finit\u{a} a lui $\Sigma$. Atunci $$\Sigma_0\se \Gamma \cup \Delta \cup \{\exists^{\geq n_1},
\ldots, \exists^{\geq n_k}\} \quad \text{pentru un } k\in\N. $$

Fie $m:=\max\{n_1,\ldots, n_k\}$. Conform (*), $\Gamma$ are un model finit $\cA$ \ai $|A|\geq m$.  Deoarece $\cA\in \cM$, avem conform (**), c\u a  $\cA\models  \Delta$. 
De asemenea, $\cA\models \exists^{\geq n_i}$ pentru orice $i=1,\ldots, k$. 
Prin urmare, $\cA\models \Gamma \cup \Delta \cup  \{\exists^{\geq n_1},
\ldots, \exists^{\geq n_k}\}$, de unde rezult\u a c\u a $\cA\models \Sigma_0$. A\c sadar, $\Sigma_0$ este satisfiabil\u a.


Aplic\^{a}nd Teorema de compacitate, rezult\u{a} c\u{a}   $\Sigma$ are un model $\cB$. 

Deoarece $\cB\models \Delta$, avem, conform (**), c\u a $\cB\in \cM$. \^In particular,   $\cB$ este finit\u{a}. 

Deoarece $\cB\models \{\exists^{\geq n}\mid n\geq 1\}$,  rezult\u{a} c\u{a} $\cB$ este infinit\u{a}.
 
 Am ob\c{t}inut o contradic\c{t}ie. 
}

\section*{Partea II. Probleme de laborator cu rezolvare clasic\u{a}}
\probex[1 punct]
    Defini\c{t}i un predicat \texttt{expand\_intervals/2} care, pentru o list\u{a} de perechi de numere 
    naturale, calculeaz\u{a} lista R format\u{a} \^{i}n felul urm\u{a}tor:
    Pentru orice \texttt{i} mai mic dec\^{a}t lungimea lui \texttt{L},
    dac\u{a} \texttt{(N, M)} este perechea de pe pozi\c{t}ia \texttt{i} din \texttt{L},
    atunci pe pozi\c{t}ia \texttt{i} din \texttt{R},
    se va afla lista tuturor numerelor mai mari sau egale ca \texttt{N} \c{s}i mai mici sau egale 
    cu \texttt{M}.

    \textbf{Exemplu:}
\begin{verbatim}
?- expand_intervals([(1, 3), [5, 5], [5, 3] [2, 6]], R).
R = [[1, 2, 3], [5], [], [2, 3, 4, 5, 6]]. 
\end{verbatim}
\textbf{Rezolvare: }
\begin{verbatim}
expand(N, M, []) :- N > M.
expand(N, M, [N|Rest]) :-
    N =< M,
    N1 is N + 1,
    expand(N1, M, Rest).

expand_intervals([], []).
expand_intervals([(N,M)|T], [L|Rest]) :-
    expand(N, M, L),
    expand_intervals(T, Rest).
\end{verbatim}


\probex[1 punct]
Consider\u{a}m \^{i}n continuare reprezentarea formulelor logicii propozi\c{t}ionale folosit\u{a} 
    \^{i}n laboratorul 5.
   Scrie\c{t}i un predicat \texttt{assoc\_and} care, 
    primind ca argument dou\u{a} formule \texttt{Phi}, \texttt{Psi}, 
    este adev\u{a}rat 
    daca \c{s}i numai dac\u{a} 
    \texttt{Phi} \c{s}i \texttt{Psi} sunt formate doar din conjunc\c{t}ii \c{s}i variabile
    \c{s}i, \^{i}n plus,
    \texttt{Psi} se poate ob\c{t}ine din \texttt{Phi}
    prin reasocierea parantezelor din ea (dar p\u{a}str\^{a}nd ordinea \^{i}n care apar variabilele \^{i}n conjunc\c{t}ii). 

    \textbf{Exemplu:}
\begin{verbatim}
?- assoc_and(si(si(a, b), c), si(a, si(b, c))). 
true 
?- assoc_and(si(si(a, b), si(c, d)), si(a, si(b, si(c, d)))).
true 
?- assoc_and(si(a, b), si(b, a)).
false 
?- assoc_and(a, a).
true 
?- assoc_and(sau(a, sau(b, c)), sau(sau(a, b), c)).
false 
\end{verbatim}
\textbf{Rezolvare:}
\begin{verbatim}
flatten_and(si(A, B), Vars) :-
    flatten_and(A, VA),
    flatten_and(B, VB),
    append(VA, VB, Vars).

flatten_and(X, [X]) :- atom(X).

assoc_and(Phi, Psi) :-
    flatten_and(Phi, Vars1),
    flatten_and(Psi, Vars2),
    Vars1 = Vars2.
\end{verbatim}

\probex[1 punct]
Definiți un predicat \texttt{dropN/3}, astfel încât, pentru orice liste \texttt{L}, \texttt{R} și număr natural \texttt{N},
\texttt{dropN(L, R, N)} este adevărat dacă și numai dacă \texttt{R} este lista care rezultă din eliminarea ultimelor
\texttt{N} elemente ale lui \texttt{L}. Predicatul va fi fals în cazul în care \texttt{N} este mai mare decât lungimea lui \texttt{L}.

\textbf{Exemplu:}
\begin{verbatim}
?- dropN([a, b, c, b], R, 2).
R = [a, b]
?- dropN([a, b, c], R, 5).
false.
\end{verbatim}

\textbf{Rezolvare:}
\begin{verbatim}
dropN(L, R, N) :- append(R, L1, L), length(L1, N).
\end{verbatim}


\section*{Partea III. Probleme de tip gril\u a}


\probex[0,5 puncte; 1 r\u{a}spuns corect] Fie urm\u{a}toarea mul\c{t}ime de formule 
\begin{align*}
\Gamma = \{ (v_i \to v_{i + 1}) \wedge (v_{i + 1} \to v_i) \mid i \geq 1 \}. 
\end{align*}
Care dintre urm\u{a}toarele afirma\c{t}ii este adev\u{a}rat\u{a}? 

$\Box$ A: $\Gamma$ are exact 2 modele. 

$\boxtimes$ B: $\Gamma$ are exact 4 modele. 

$\Box$ C: $Mod(\Gamma \cup \{v_1\}) \subseteq Mod(\Gamma \cup \{v_0\})$.

$\Box$ D: $\Gamma$ are o infinitate de modele. 

$\Box$ E : $\Gamma \cup \{ \neg v_1 \to v_2 \}$ este nesatisfiabil\u{a}.

\solutiong{
    Observ\u{a}m c\u{a} pentru orice evaluare $e$, avem c\u{a} 
    \begin{align*}
        &\iff e \in Mod(\Gamma) \\ 
        &\iff \text{ pentru orice $i \geq 1$, } e^+(v_i \leftrightarrow v_{i + 1}) = 1 \\
        &\iff \text{ pentru orice $i \geq 1$, } e(v_i) = e(v_{i + 1}) \\
        &\iff \text{ pentru orice } i \geq 1, e(v_i) = 0 \text{ sau pentru orice } i \geq 1, e(v_i) = 1
    \end{align*}
    A\c{s}adar, modelele lui $\Gamma$ sunt exact evalu\u{a}rile $e_0, e_1, e_2, e_3 : V \to \{0, 1\}$ 
    definite prin
    \begin{align*}
        \text{pentru orice $i \in \N$},\ e_0(v_i) = 0 \\ 
        \text{pentru orice $i \in \N$},\ e_1(v_i) = 1 \\ 
        e_2(v_0) = 0 \ \text{\c{s}i pentru orice $i \geq 1$},\ e_2(v_i) = 1 \\ 
        e_3(v_0) = 1 \ \text{\c{s}i pentru orice $i \geq 1$},\ e_3(v_i) = 0.
    \end{align*}
    Deci B este adev\u{a}rat\u{a}, iar A, D sunt false. \\ 
    C: Avem c\u{a} $e_2$ definit mai sus este model pentru $\Gamma \cup \{v_1\}$,
    dar nu \c{s}i pentru $\Gamma \cup \{v_0\}$. \\
    E: $e_1$ este model pentru $\Gamma \cup \{ \neg v_1 \to v_2 \}$.
}

\probex[0,5 puncte; 1 r\u{a}spuns corect] Pentru orice mul\c{t}ime de formule propozi\c{t}ionale $\Gamma \subseteq Form$,  
definim $\overline{\Gamma} = \{\neg \varphi \mid \varphi \in \Gamma\}$. 
Care dintre urm\u{a}toarele afirma\c{t}ii este adev\u{a}rat\u{a} pentru orice $\Gamma \subseteq Form$ nevid\u{a}?

$\Box$ A: Pentru orice evaluare $e$, avem c\u{a} $e \in Mod(\overline{\Gamma})$ \ddaca $e \notin Mod(\Gamma)$.

$\Box$ B: Dac\u{a} $\Gamma$ este nesatisfiabil\u{a}, atunci $\overline{\Gamma}$ este satisfiabil\u{a}.

$\Box$ C: $Mod(\Gamma) \neq Mod(\overline{\Gamma})$.

$\Box$ D: $\Gamma \cap \overline{\Gamma}$ este nesatisfiabil\u{a}.

$\boxtimes$ E: $\Gamma \cup \overline{\Gamma}$ este nesatisfiabil\u{a}.

\solutiong{
    Contraexemplul $\Gamma = \{v_0, \neg v_0\}$ arat\u{a} c\u{a} A, B, C, D sunt false. \\ 
    E: Presupunem c\u{a} exist\u{a} $e \in Mod(\Gamma \cup \overline{\Gamma})$. 
    Deoarece $\Gamma$ este nevid\u{a}, putem alege $\varphi \in \Gamma$. 
    Vom avea c\u{a} $\neg \varphi \in \overline{\Gamma}$, 
    de unde ob\c{t}inem c\u{a} $e^+(\varphi) = e^+(\neg \varphi) = 1$, 
    ceea ce este o contradic\c{t}ie.
}

\probex [0,5 puncte; 2 r\u{a}spunsuri corecte]
Fie urm\u{a}toarea formul\u{a} propozi\c{t}ional\u{a}: 
\begin{align*}
\varphi = (v_0 \wedge v_1 \wedge v_2 \wedge v_3) \to v_4. 
\end{align*}
Care dintre urm\u{a}toarele afirma\c{t}ii sunt adev\u{a}rate? 

$\boxtimes$ A: $\varphi \sim v_0 \to (v_1 \to (v_2 \to (v_3 \to v_4)))$.

$\Box$ B: $\varphi \sim (((v_0 \to v_1) \to v_2) \to v_3) \to v_4$.

% $\Box$ $v_1 \to v_4 \models \varphi$

$\Box$ C: $\varphi \models v_1 \to v_4$.

$\Box$ D: $\varphi \models (v_0 \vee v_1 \vee v_2 \vee v_3) \to v_4$ .

% $\Box$ $(v_0 \vee v_1 \vee v_2 \vee v_3) \to v_4 \models \varphi$

$\boxtimes$ E: $\neg v_0 \vee \neg v_1 \vee \neg v_2 \vee \neg v_3 \vee v_4$ este FND \c{s}i FNC pentru $\varphi$.

\solutiong{
    A: se aplic\u{a} succesiv echivalen\c{t}a $\varphi \wedge \psi \to \chi \ \sim \ \varphi \to (\psi \to \chi)$. \\ 
    E: Avem 
    \begin{align*}
        &(v_0 \wedge v_1 \wedge v_2 \wedge v_3) \to v_4 \quad\sim\\ 
        &\neg(v_0 \wedge v_1 \wedge v_2 \wedge v_3) \vee v_4 \quad\sim \\ 
        &\neg v_0 \vee \neg v_1 \vee \neg v_2 \vee \neg v_3 \vee v_4,
    \end{align*}
    unde ultima formul\u{a} este at\^{a}t \^{i}n FND, c\^{a}t \c{s}i \^{i}n FNC.

    Fie $e$ o evaluare cu $e(v_1) = 1$ \c{s}i $e(v_0) = e(v_2) = e(v_3) = e(v_4) = 0$.
    Avem \\ 
    B: $e \models \varphi$, dar $e \not{\models} (((v_0 \to v_1) \to v_2) \to v_3) \to v_4$ \\ 
    C: $e \models \varphi$, dar $e \not{\models} v_1 \to v_4$ \\
    D: $e \models \varphi$, dar $e \not{\models} (v_0 \vee v_1 \vee v_2 \vee v_3) \to v_4$
    
}

\probex 
[0,5 puncte; 2 r\u{a}spunsuri corecte]
Fie $\mathcal{L}$ un limbaj de ordinul \^{i}nt\^{a}i care con\c{t}ine exact un simbol de func\c{t}ie de aritate 2,
notat $\dot{\times}$. 
Consider\u{a}m urm\u{a}toarea formul\u{a} a lui $\mathcal{L}$
\begin{align*}
\varphi = \forall x \forall y (x = y \dot{\times} y \to \forall z (x = z \dot{\times} z \to y = z)),
\end{align*} 
unde $x, y$ \c{s}i $z$ sunt variabile distincte dou\u{a} c\^{a}te dou\u{a}.
Care dintre urm\u{a}toarele afirma\c{t}ii sunt adev\u{a}rate? 

$\boxtimes$ A: $(\N, \cdot) \models \varphi$. 

$\Box$ B: $(\Q, \cdot) \models \varphi$.

$\boxtimes$ C: $\varphi \vDashv \forall x \forall y \forall z (x = y \dot{\times} y \to (x = z \dot{\times} z \to y = z))$.

$\Box$ D: $\varphi \vDashv \forall x \forall y \exists z (x = y \dot{\times} y \to (x = z \dot{\times} z \to y = z))$.

$\Box$ E: $\varphi \vDashv \forall x (\exists y (x = y \dot{\times} y) \to \forall z (x = z \dot{\times} z \to y = z))$.

\solutiong{
    Observ\u{a}m c\u{a} $\varphi$ este satisf\u{a}cut \^{i}n 
    $(\N, \cdot)$ (sau $(\Q, \cdot)$) dac\u{a} \c{s}i numai dac\u{a} 
    pentru orice num\u{a}r din $\N$ (resp. din $\Q$), dac\u{a} acesta are o 
    `r\u{a}d\u{a}cin\u{a} p\u{a}trat\u{a}', aceasta este unic\u{a}.  
    Afirma\c{t}ia este adev\u{a}rat\u{a} \^{i}n $\N$, dar nu \c{s}i \^{i}n 
    $\Q$ (e.g. $1 \cdot 1 = (-1) \cdot (-1) = 1$).
    Formula de la C este o form\u{a} normal\u{a} prenex pentru $\varphi$.
}

\probex [0,5 puncte; 2 r\u{a}spunsuri corecte]
Fie $\mathcal{L}$ un limbaj de ordinul \^{i}nt\^{a}i con\c{t}in\^{a}nd 
un simbol de rela\c{t}ie binar\u{a} $P$ \c{s}i un simbol de rela\c{t}ie unar\u{a} 
$Q$.
Fie urm\u{a}toarea formul\u{a} a lui $\mathcal{L}$
\begin{align*}
\varphi = \forall x \exists y \forall u \exists z 
(P(x, y) \to (P(u, u) \to \neg Q(z))), 
\end{align*}
unde $x, y, u$ \c{s}i $z$ sunt variable distincte dou\u{a} c\^{a}te dou\u{a}.
Care dintre urm\u{a}toarele afirma\c{t}ii sunt adev\u{a}rate? 

$\boxtimes$ A: $\varphi$ este o form\u{a} normal\u{a} prenex pentru formula 
\begin{align*}
\exists x \forall y P(x, y) \to (\neg \forall x \neg P(x, x) \to \neg \forall z Q(z)).
\end{align*}

$\boxtimes$ B: $\forall x \forall u (P(x, g(x)) \to (P(u, u) \to \neg Q(h(x, u))))$ este o form\u{a} normal\u{a} Skolem pentru $\varphi$,
unde $g$ \c{s}i $h$ sunt simboluri noi de func\c{t}ie de aritate $1$, respectiv $2$.

$\Box$ C: $\varphi$ este o form\u{a} normal\u{a} prenex pentru formula 
\begin{align*}
\exists x \forall y P(x, y) \to (\exists x P(x, x) \to \forall z Q(z)).
\end{align*}

$\Box$ D: $\forall x \forall u (P(x, g(x, u)) \to (P(u, u) \to \neg Q(h(x, u))))$ este o form\u{a} normal\u{a} Skolem pentru $\varphi$,
unde $g$ \c{s}i $h$ sunt simboluri noi de func\c{t}ie de aritate $2$. 

$\Box$ E: $\varphi$ este at\^{a}t \^{i}n form\u{a} normal\u{a} prenex, c\^{a}t \c{s}i \^{i}n form\u{a} normal\u{a} Skolem.

\solutiong{
Varianta A este corect\u{a} pentru forma prenex: 
\begin{align*}
    \exists x \forall y P(x, y) \to (\neg \forall x \neg P(x, x) \to \neg \forall z Q(z)) \vDashv \\ 
    \exists x \forall y P(x, y) \to (\exists x P(x, x) \to \exists z \neg Q(z)) \vDashv \\ 
    \exists x \forall y P(x, y) \to (\exists u P(u, u) \to \exists z \neg Q(z)) \vDashv \\ 
    \exists x \forall y P(x, y) \to \forall u \exists z (P(u, u) \to \neg Q(z)) \vDashv \\ 
    \forall x \exists y (P(x, y) \to \forall u \exists z (P(u, u) \to \neg Q(z))) \vDashv \\ 
    \forall x \exists y \forall u \exists z (P(x, y) \to (P(u, u) \to \neg Q(z))) \\ 
\end{align*}
}

\probex [0,5 puncte; 1 r\u{a}spuns corect] Consider\u{a}m urm\u atoarea mul\c time de clauze:
\begin{align*}
    \mathcal{S} = \{ \{v_0, v_1, \neg v_2\}, \{v_0, \neg v_1, v_2\}, \{\neg v_0, v_1, v_2\}, \{\neg v_0, \neg v_2\} \}
\end{align*}
Care dintre urm\u{a}toarele afirma\c{t}ii este adev\u{a}rat\u{a}?

$\boxtimes$ A: $\mathcal{S}$ este satisfiabil\u{a}. 

$\Box$ B: $\mathcal{S}$ este nesatisfiabil\u{a}. 

$\Box$ C: $\{v_1\}$ este rezolvent al dou\u{a} clauze din $\mathcal{S}$. 

$\Box$ D: $\{v_0\}$ este rezolvent al dou\u{a} clauze din $\mathcal{S}$. 


$\Box$ E: Rul\^{a}nd algoritmul Davis-Putnam pe mul\c{t}imea $\mathcal{S}$,
vom ob\c{t}ine $\square \in \mathcal{S}_{N + 1}$,
unde $N$ este num\u{a}rul de pa\c{s}i dup\u{a} care algoritmul se termin\u{a}. 


\solutiong{
    Se observ\u{a} c\u{a} o evaluare $e$ cu 
    $e(v_0) = e(v_1) = 1$ \c{s}i $e(v_2) = 0$ satisface $\mathcal{S}$. 
    Rezult\u{a} c\u{a} A este adev\u{a}rat\u{a} \c{s}i B, E sunt false.
    Se observ\u{a} c\u{a} $\{v_0\}$ \c{s}i $\{v_1\}$ nu se pot ob\c{t}ine ca rezolven\c{t}i din dou\u{a} clauze din $\mathcal{S}$.
}


\probex [0,5 puncte; 1 r\u{a}spuns corect]
Consider\u{a}m urm\u{a}torul program Prolog.
\begin{verbatim}
p(Xs, [0|Rest]) :-
    p1(Xs, 0, 1, Rest).

p1([], _, _, []).
p1([X|Xs], S, P, [P2|Rest]) :-
    S2 is S + X,
    P2 is P * X,
    p2(Xs, S2, P2, Rest). 

p2([], _, _, []).
p2([X|Xs], S, P, [S2|Rest]) :-
    S2 is S + X,
    P2 is P * X,
    p1(Xs, S2, P2, Rest). 
\end{verbatim}
Care dintre urm\u{a}toarele afirma\c{t}ii este adev\u{a}rat\u{a}? \\ 
$\Box$ R\u{a}spunsul interogarii \texttt{?- p([2, 2, 2, 2, 2, 2], R)}. este \texttt{R = [0, 2, 4, 8, 16, 32]}.

$\Box$ R\u{a}spunsul interogarii \texttt{?- p([1, 1, 1, 1]).} este \texttt{R = [0, 1, 2, 3]}

$\Box$ R\u{a}spunsul interogarii \texttt{?- p([], []).} este \texttt{true}.

$\Box$ R\u{a}spunsul interogarii \texttt{?- p([], R).} este \texttt{false}. 

$\boxtimes$ R\u{a}spunsul interogarii \texttt{?- p([2, 2, 2, 2, 2, 2], R).} este \texttt{R = [0, 2, 4, 8, 8, 32, 12]}

\probex [0,5 puncte; 2 r\u{a}spunsuri corecte]

Consider\u{a}m urm\u{a}torul program Prolog: 
\begin{verbatim}
p([], []).

p([_|T], S) :- p(T, S).

p([X|T], [X|S]) :- p_from(T, X, S).

p_from([], _, []).

p_from([_|T], Last, S) :- p_from(T, Last, S).

p_from([X|T], Last, [X|S]) :- X > Last, p_from(T, X, S).
\end{verbatim}
Care dintre urm\u{a}toarele afirma\c{t}ii sunt adev\u{a}rate? \\ 
$\boxtimes$ R\u{a}spunsul interogarii \texttt{p([4,6,8,3,4,78,5,2,6], [3, 4, 5, 6])} este \texttt{true}.

$\Box$ R\u{a}spunsul interogarii \texttt{p([4,6,8,3,4,78,5,2,6], [4, 3])} este \texttt{true}.

$\Box$ Pentru orice list\u{a} \texttt{L}, exist\u{a} o singur\u{a} valoare posibil\u{a} pentru \texttt{R} 
astfel \^{i}nc\^{a}t interogarea \texttt{?- p(L, R)} s\u{a} aib\u{a} r\u{a}spunsul \texttt{true}. 

$\boxtimes$ R\u{a}spunsul interogarii \texttt{p([6, 5, 4, 3, 2, 1], [])} este \texttt{true}. 

$\Box$ Nu exist\u{a} nicio valoare pentru \texttt{R} astfel \^{i}nc\^{a}t 
r\u{a}spunsul interog\u{a}rii \texttt{p([6, 5, 4, 3, 2, 1], R)} s\u{a} fie \texttt{true}


\end{document}